\section{Introduction}\label{sec:intro}

Under one controlled protocol with unbiased test data, the gains reported for exposure-bias corrections of graph
recommenders depend strongly on evaluation choices, and third-party corrections add little over an equally tuned
backbone.

Graph-based recommenders represent users and items as nodes and score a user-item pair by propagating embeddings over
the user-item graph. LightGCN \citep{he2020lightgcn} is a standard backbone of this family, and it is usually trained on
implicit feedback logs. In such a log an interaction is recorded only when the item was shown and the user responded.
The logged graph therefore reflects the exposure policy as well as user taste \citep{marlin2009collaborative,chen2023bias},
and a model fitted to the log treats every absence as informative. The pairwise loss that trains LightGCN, Bayesian
personalised ranking \citep{rendle2009bpr}, draws its negatives uniformly from the items the user has no logged
interaction with. Under exposure bias most of these items were never shown. Corrections act on the loss
\citep{schnabel2016,wang2019drjl,guo2021mrdr}, on the latent exposure or on the likelihood weights
\citep{liang2016exposure,saito2020unbiased}, and, for graph models, on aggregation or adjacency
\citep{kim2022navip,zhao2022radjnorm,islam2026ppd}.

This paper starts from the two nuisance functions that a DR estimator already fits: an exposure propensity $p_{ui}$ and
an imputed preference $\Yhat_{ui}$. Under an independence assumption between exposure and preference given covariates,
Bayes' rule gives the probability that a pair absent from the log is a positive,
\begin{equation}\label{eq:Qintro}
Q_{ui}=\Pr\bigl(Y_{ui}=1\mid O_{ui}Y_{ui}=0,\,x_{ui}\bigr)
=\frac{\Yhat_{ui}\,(1-p_{ui})}{1-p_{ui}\,\Yhat_{ui}}.
\end{equation}
The candidate models use $Q$ in two places inside an otherwise unchanged LightGCN. Each sampled BPR negative is
down-weighted by $1-Q_{uj}$, and the standardised logit of the imputation is added to the score as a fixed offset whose
weight is chosen on validation data. Both nuisance functions are fitted on the biased log before training. The one
exception to the rule that training uses no uniformly collected data is Coat, whose shipped propensities were estimated
by the dataset authors from a uniform sample. We call the full model QP-LightGCN and evaluate it with its variants as one
family.

The gains reported for these corrections are hard to compare, for three reasons. First, the studies use uniformly
collected data in different ways: some train on it, some select configurations on it, and some use it only for testing
(Section~\ref{sec:related}). Second, tuning budgets differ, so a method given a larger search space can obtain a higher
test score through selection alone. Third, single runs are common, and the ordering of close methods can depend on the
random seed. Our own runs illustrate the third point. In the main run, QP-LightGCN is above iALS on Coat by
$+20.02\times10^{-3}$ [19.74], which is nominally significant after Holm correction over the eight baseline comparisons
($p_\mathrm{Holm}=0.047$). When the selected configurations are retrained with seeds 1 to 3, the difference is
$+3.39\times10^{-3}$ [20.70], with $p_\mathrm{Holm}=0.75$ over the two baseline comparisons of that check. These seed
comparisons are descriptive, because the two baselines of each dataset were chosen on test data.

The study asks four research questions, which Table~\ref{tab:rq} answers and locates.

\begin{table}[htbp]\centering\footnotesize
\caption{Research questions, short answers and the sections that answer them.}\label{tab:rq}
\begin{tabular}{p{4.6cm}p{7.2cm}l}
\toprule
Question & Short answer & Section \\
\midrule
RQ1: How much do evaluation choices change the conclusions? & A lot. Selection on the test split inflates scores and can reverse an order. Many pairwise signs depend on the search budget, the seed or the metric. & Section~\ref{sec:res-eval} \\
RQ2: How much do graph-specific corrections add over an equally tuned backbone? & Little, on Coat, Yahoo! R3 and KuaiRand-Pure. On KuaiRec the gains are larger, and a pointwise loss without debiasing gives a larger difference of means. & Section~\ref{sec:res-debias} \\
RQ3: Which ingredient carries the gains that remain? & An additive prior fitted once on the exposed log. Its inverse-propensity weights are not what matters in our runs, and the mechanism is not identified. & Section~\ref{sec:res-ingredient} \\
RQ4: Does exact unbiasedness of the propagation matter for ranking? & Not on the three real logs on which it was run. & Section~\ref{sec:res-walk} \\
\bottomrule
\end{tabular}
\end{table}

The paper makes four contributions. None of them is a new estimator.
\begin{enumerate}
\item \emph{Protocol and its sensitivity (RQ1).} Four public benchmarks with unbiased test data, uniform data used only for
validation and testing, per-split selection on validation data, per-user paired tests with Holm correction, and
retraining with seeds. Selecting on the test split inflates Coat scores by up to $0.030$ nDCG@5.
\item \emph{Third-party debiasers add little over an equally tuned backbone (RQ2).} On Coat, Yahoo! R3 and KuaiRand-Pure
the largest point-estimate gain of a third-party debiaser over LightGCN is $+4.10\times10^{-3}$ nDCG, and none of these
gains survives a correction for overlapping splits on those datasets.
\item \emph{Prior anatomy (RQ3).} The gains over the backbone that the additive prior provides are examined with controls.
Removing its inverse-propensity weights does not reduce them, a shuffled item term removes the KuaiRec gain, and
popularity does not reproduce it. The mechanism is not identified.
\item \emph{A negative result on walk correction (RQ4).} The walk correction of DRUP removes the repeated-edge bias of
multi-hop DR propagation. It does not improve ranking on Coat, Yahoo! R3 or KuaiRec.
\end{enumerate}

The QP-LightGCN family and its refinements are members of the evaluated family (\ref{app:refine}), not
contributions in their own right. We do not claim that debiasing is useless in general.

The scope is four public benchmarks, offline evaluation, top-$K$ nDCG at the cut-off used for each dataset (5 for Coat
and Yahoo! R3, 10 for KuaiRand-Pure, 20 for KuaiRec), and implicit-feedback learners trained on the biased log. The
walk-correction claim concerns Coat, Yahoo! R3 and KuaiRec, the three real logs on which DR and DRUP were run. Coat has
225 evaluated users, so its power is low. For Coat we report the differences that the data exclude, and we do not report
equality. We do not claim that AdvInfoNCE, NAVIP, r-AdjNorm, SimGCL or MACR fail in general. Our implementations, tuned
with the stated budgets under this protocol, do not beat the backbone by more than the amounts reported. We do not give a
mechanism for why the prior helps beyond the controls of Section~\ref{sec:res-ingredient}, and the reading that the
prior acts as a debiased item-quality score is a hypothesis.
